\documentclass[10pt,a4paper]{amsart}
\usepackage[margin=19mm]{geometry}
\usepackage{amsmath,amssymb,mathtools}
\usepackage[scr=rsfs,cal=euler]{mathalfa}
\usepackage{xfrac}
\usepackage[unicode,hidelinks,bookmarks=false]{hyperref}
\newcommand{\ie}{i.e.\ }
\newcommand{\cf}{cf.\ }
\newcommand{\resp}{respectively\ }
\newcommand{\Subsection}{\subsection}
\providecommand{\nobreakdash}{\nobreak\hbox{-}}
\setlength{\emergencystretch}{2em}
\multlinegap0pt
\usepackage{amssymb}
\usepackage{enumerate}
\usepackage{tikz}
\usepackage{dsfont}
\usepackage{xcolor}
\newtheorem{proposition}{Proposition}
\newcommand{\AZ}{\textcolor{blue}}
\def\SO{{\mathcal O}}
\newcommand{\be}{\begin{eqnarray}}
\newcommand{\bea}{\begin{eqnarray*}}
\newcommand{\ee}{\end{eqnarray}}
\newcommand{\eea}{\end{eqnarray*}}
\def\ml{\lambda}
\title{quantisation in high dimensional sets}
\author{Luc Pronzato, Anatoly Zhigljavsky}
\date{Source: arXiv:2605.12568v1 (CC BY 4.0)}
\begin{document}
\maketitle

\noindent\emph{Source and licence.} quantisation in high dimensional sets. Version arXiv:2605.12568v1 (CC BY 4.0), distributed by arXiv under the Creative Commons Attribution 4.0 International licence, http://creativecommons.org/licenses/by/4.0/, as recorded in the arXiv licence metadata for this exact version and shown on the abstract page https://arxiv.org/abs/2605.12568v1. Full licence notice: \texttt{licenses/arxiv-2605.12568-CC-BY-4.0.txt}.\par\smallskip

\noindent\textit{Editorial scope.} Counted unit: the article's proposition th:kappa\_n (source0.tex lines 906--917). The article's complete original proof of it is delivered with it (source0.tex lines 919--966, one \texttt{proof} environment of 48 lines). No mathematics is written, repaired or re-derived: the statement and the proof are the article's own lines, in the article's own notation, and no retained line is substituted or re-flowed.  No further source-local context is needed: every cross-reference of every retained line resolves inside the two delivered spans.  Nothing was omitted from any retained span, so the package carries no omission record.  No retained line contains a citation, so the wrapper carries no bibliography and no bibliography entry.  No retained line contains a figure environment, an image-inclusion command or drawing code, so no image is delivered and the package declares no asset dependency.  Editorial formatting only, in addition: an AMS \texttt{amsart} preamble that restates the article's own declarations and macros used by the retained lines, namely the environments \texttt{proposition} on separate counters, together with the standard packages those lines need.  The retained environments print in this document as Proposition 1 in order of appearance, whereas the article prints the same items with the numbering of its own full document.  The complete native source of the article as distributed in the arXiv e-print for this exact version is bundled unchanged as \texttt{sources/200427/source0.tex}.
\begin{proposition} \label{th:kappa_n}
For fixed $d\geq 3$, the $1/n$ quantile $\kappa_{n,d}$ of  $\beta_{\delta,\delta}$ with $\delta=(d-1)/2$ tends to zero as $n$ tends to infinity. When $d \to \infty$ and $n=n(d)$ grows with $d$, we have the following three cases:\\
\noindent{\textit(i)} if $n^{1/d}\to\infty$, then $\lim_{d\to\infty} \kappa_{n,d} = 0$;\\
\noindent{\textit(ii)} if $n=C\ml^d\,(1+o(1))$ with $\ml>1$ and $C>0$, then
\be \label{asymptotic-kappa} \lim_{d\to\infty} \kappa_{n,d} = \frac12 \left[1-\sqrt{1-1/\ml^2}\right]\, ;
\ee
%\be\label{bounds-kappa-caseii}
%\frac{1}{4\,\ml^2} \leq \liminf_{d\to\infty} \kappa_{n,d} \leq \limsup_{d\to\infty} \kappa_{n,d} \leq \frac12\,\left[1-\sqrt{1-1/\ml^2}\right]\,;
%\ee
\noindent{\textit(iii)} %if $n=d^\mt\,(1+o(1))$ with $\mt>0$,
if $(\log n)/d \to 0$, then $\lim_{d\to\infty} \kappa_{n,d} = \frac12$.
\end{proposition}

\begin{proof}
Since $x^{\delta-1}(1-t)^\delta\leq x^{\delta-1}(1-t)^{\delta-1}\leq x^{\delta-1}(1-x)^{\delta-1}$ for $\delta\geq 1$ and all $t\in[0,1)$ and $x\in[0,t]$, we have
\be\label{It-ineq1}
\frac{t^\delta (1-t)^\delta}{\delta B(\delta,\delta)} \leq I_t(\delta,\delta) \,. %\leq \frac{t^\delta}{\delta B(\delta,\delta)}\,.
\ee
The derivation of an exploitable upper bound is more delicate.
Assume that $d \geq 3$ and let
$\delta'=\lfloor \delta \rfloor$, the largest integer smaller than or equal to $\delta$.
%Assume that $d\geq 3$ and let $\delta'$ be the largest integer $m$ such that $2\,m+1 \leq d$.
As $n\geq 2$, we are only interested in values of $t$ less than $1/2$. Since $\delta'\leq \delta$, we have $I_t(\delta,\delta) \leq I_t(\delta',\delta')$, which
%For any $n\geq 2$, we have $\kappa_{n,d}\leq \kappa_{n,2\,\delta'+1}$, where $\kappa_{n,2\,\delta'+1}$ is the solution of $I_t(\delta',\delta')=1/n$ with respect to $t$ and where the c.d.f.\ $I_t(\delta',\delta')$
can be calculated explicitly by successive integration by parts. Direct calculations yield the following special case of the well-known relation between the c.d.f.\ of the beta and binomial distributions:
\be\label{It1}
%I_t(\delta',\delta') = \frac{1}{\delta' B(\delta',\delta')} \, \sum_{j=1}^{\delta'} A_j\, t^{\delta'+j-1}(1-t)^{\delta'-j} \,,
I_t(\delta',\delta') &=& \sum_{k=0}^{\delta'-1} \binom{2\delta'-1}{k} t^{2\delta'-1-k}(1-t)^k \,, \\
&=& t^{\delta'}(1-t)^{\delta'-1} \sum_{k=0}^{\delta'-1} \binom{2\delta'-1}{k} \left(\frac{t}{1-t}\right)^{\delta'-1-k} \,.
\ee
%where
%\bea
%A_j= \frac{\delta'(\delta'-1)(\delta'-2)\cdots(\delta'-j+1)}{\delta'(\delta'+1)(\delta'+2)\cdots(\delta'+j-1)} = \frac{\delta'!(\delta'-1)!}{(\delta'-j)!(\delta'+j-1)!} \,.
%\eea
%We thus obtain
%\bea
%I_t(\delta',\delta') = \frac{t^{\delta'}(1-t)^{\delta'-1}}{\delta' %B(\delta',\delta')} \, \sum_{j=1}^{\delta'} A_j\, t^{j-1}(1-t)^{1-j} \leq \frac{t^{\delta'}(1-t)^{\delta'-1}}{\delta' B(\delta',\delta')} \sum_{j=1}^{\delta'} A_j \,,
%\eea
%where we used the property $t/(1-t) \leq 1$ for $t\leq 1/2$. Since $I_{1/2}(\delta',\delta')=1/2$, \eqref{It1} gives $\sum_{j=1}^{\delta'} A_j=4^{\delta'-1}\,\delta' B(\delta',\delta')$, so that
Using the property $t/(1-t) \leq 1$ for $t\leq 1/2$, with $\sum_{k=0}^{\delta'-1} \binom{2\delta'-1}{k}=4^{\delta'-1}$, we thus obtain
\be \label{upper_bound}
I_t(\delta,\delta) \leq I_t(\delta',\delta') \leq t^{\delta'}[4\,(1-t)]^{\delta'-1} \leq \frac12\,[4\,t(1-t)]^{\delta'-1}\,, \ t\leq 1/2\,.
\ee
%\footnote{\AZ{New: you simply removed one $t$ by using $t\leq 1$ but you could have used $t\leq 1/2$ and gained the multiplier $1/2$ in \eqref{upper_bound}. The lower bound in \eqref{bounds-kappa} would slightly improve. \\Also, to get the intermediate inequality you may perhaps use the inequality $1/(1-t)\leq 2 $ rather than $t/(1-t)\leq 1 $. In this way, I think we would  keep the correct value of the shape parameter in the intermediate inequality if $\delta=\delta'$. } }

Together with \eqref{It-ineq1}, this implies, for $n\geq 2$ and $d\geq 3$,
\be\label{bounds-kappa}
\frac12 \left[1-\sqrt{1- (2/n)^{1/(\delta'-1)}}\right]  \leq \kappa_{n,d} \leq \frac12 \left[1-\sqrt{(1-4\,c_d\,n^{-1/\delta})_+}\right] \,,
%c_d\,n^{-1/\delta} \leq \kappa_{n,d} \leq C_{n,d} = \left\{ \begin{array}{ll}
%\frac12 \left[1-\sqrt{1-4\,c_d\,n^{-1/\delta} }\right] & \mbox{ if } n > (4\,c_d)^\delta\,, \\
%\frac12 & \mbox{ otherwise}\,, \end{array}\right.
\ee
where $(x)_+=\max\{x,0\}$ and $c_d=[\delta\,B(\delta,\delta)]^{1/\delta} = 1/4+ (\log d)/(4\,d)+\SO(1/d)$, $d\to\infty$ (with, moreover, $(4\,c_d)^\delta<d+2$ for all $d \geq 2$, implying that the right-hand side of \eqref{bounds-kappa} is strictly smaller than  1/2 for $n\geq d+2$).

When $d$ is fixed and $n\to\infty$, or $n^{1/d}\to\infty$ as $d\to\infty$, $n^{-1/d}\to 0$ implying that the right-hand side of \eqref{bounds-kappa} tends to zero. Therefore, $\kappa_{n,d}\to 0$.

When $n=C\ml^d\,(1+o(1))$ with $\ml>1$ and $C>0$ (case (\textit{ii})), by taking the limits on left and right hand sides of \eqref{bounds-kappa}, as $\lim_{d\to\infty} c_d=1/4$, we obtain $\lim_{d\to\infty} \kappa_{n,d} = (1/2)\,\left[1-\sqrt{1-1/\ml^2}\right]$.

%When $n=n(\mt,d)=d^\mt\,(1+o(1))$ with $\mt>0$
When $\log n = o(d)$, (case (\textit{iii})), $n^{-1/d}\to 1$ and the left-hand side of \eqref{bounds-kappa} tends to 1/2.
\end{proof}

\end{document}
